\section*{Limitations}
\label{sec:limitations}

\paragraph{One training run per $K$.}
Every effect size reported here is computed across the 96 patient personas \emph{within} a single
training run. There are no training-seed replicates, so run-to-run training variance is unmeasured,
and we cannot separate ``look-ahead helps'' from ``this look-ahead run went well.'' The $\Kz$
run's decline over its last two iterations and its late praise spike are likewise events of one
run, which is why \S\ref{sec:reward} reports the contrast against that run's best checkpoint as
well as its last, and why \S\ref{sec:therapist}--\ref{sec:patient} report at which of the ten
iterations each process contrast is significant rather than the endpoint alone. We regard this as the principal limitation; it is distinct from, and not bounded
by, the evaluation noise discussed next.

\paragraph{Evaluation draws.}
Therapist decoding is unseeded and each trained model state was evaluated on one
96-conversation sample (the Base on two, which we pool). At the trained $\Kf$ endpoint the re-draw
of \S\ref{sec:reward} differs from the original by at most $|\dz|\,0.174$ across 9 instruments and
2 judges, nothing significant. The remaining states are single draws, and a replicate bounds evaluation noise only,
never the training variance above.

\paragraph{Matched iterations are not matched cost.}
An iteration is 96 conversations and $G{=}8$ scored candidates per prompt in both runs, so oracle
scoring calls ($302{,}541$ for $\Kz$ vs.\ $289{,}983$ for $\Kf$ over ten iterations) and
optimizer steps ($1{,}128$ vs.\ $1{,}070$) are approximately matched by construction. But each
$\Kf$ reward requires a five-turn simulated continuation, about $393$k patient-simulator calls
over ten iterations that the $\Kz$ run never makes, and roughly doubles per-step wall-clock (median
$1.92\times$ over the settled iterations 3--10; $27.9$ vs.\ $51.2$ GPU-hours for the two runs).
More compute per iteration is built into the lever, and we disclose the asymmetry rather than
equalise it. Read at matched GPU-hours instead, with each run at its best checkpoint within the
budget under the selecting judge, the comparison is budget-dependent: at about 13 GPU-hours the
$\Kz$ run is ahead on \QQ\ ($\dz\,-0.74$ under the training oracle, $-0.78$ held out), at about 23
GPU-hours the two are level under the training oracle ($\dz\,0.07$, n.s.) and $\Kf$ leads under
the held-out judge ($\dz\,0.33$, $p_{\text{holm}} = .012$), and beyond the $\Kz$ run's total
budget the comparison is the best-checkpoint one of \S\ref{sec:reward}. GPU-hours are
reconstructed from artifact timestamps (Appendix~\ref{app:repro}), so this reading is coarser
than the iteration-matched one.

\paragraph{$K \in \{0, 5\}$ only.}
There is no dose--response evidence between the two levels and nothing above 5, so nothing here
says the effect is monotone in $K$, that 5 is optimal, or that a shorter horizon would not capture
most of the benefit. An intermediate $K$ is the most valuable single follow-up.

\paragraph{One optimiser, one regime.}
\citet{baruch2025pto} introduced the lever with preference trees and DPO in an easier regime of
this task (a 7B policy, more cooperative patients, a weaker judge); other optimisers and regimes
are outside this paper's scope. Evidence that would overturn the paper's reading: a second
\emph{training} run reversing the ordering, an intermediate $K$ showing a threshold rather than a
horizon, or a human MI coder disagreeing with \S\ref{sec:therapist}--\ref{sec:patient}.

\paragraph{The look-ahead reward is also a reward for continuing.}
A simulated session can end inside the look-ahead window, most often because the patient closes
it, and such a candidate gives the oracle a shorter transcript. In the rollout audit
(Appendix~\ref{app:tables}, Figure~\ref{fig:tails}) $82\%$ of logged $\Kf$ rollouts ran the full
five turns; early-ending candidates scored at or below their group mean, by up to $0.09$ points
depending on how the rollout ended, and were the group's best less often than chance. The lever was
therefore applied as specified in the large majority of cases, but it carries a mild pressure to
keep the session open that a turn-level reward does not, and the $\Kf$ policy's longer endpoint
sessions ($31.9$ vs.\ $25.2$ utterances, $\dz\,0.43$) are consistent with it. Two runs cannot separate this from the
horizon effect proper.

\paragraph{Simulation only, in sample, without human validation.}
Every conversation is between two language models, and every number is in-sample with respect to
the patient distribution: the 96 personas are the exhaustive enumeration of the persona grid, and
the self-training loop uses the same 96 for training rollouts and for evaluation at every
iteration, so generalisation to unseen patients is unmeasurable without authoring new personas.
The held-out judge decouples the grading half, but every conversation was generated against a
\texttt{gpt-4o-mini} patient, the same model as the training oracle; a fully decoupled replication
would regenerate all conversations with a patient from another model family. The simulated
patient is a known weak point in itself: given identical profiles, LLM-simulated MI patients
express themselves more uniformly positively than human patients do
\citep{hoang2026standardised}, which bears directly on a finding about praise. No human MI coder has
rated any conversation in this experiment; the oracle is highly self-repeatable (ICC 0.92--0.99 on
the two $K{=}0$ states with a repeat draw), but repeatability is not validity, and the over-praise
finding is the most robust behavioural claim here precisely because the keyword marker
corroborates it without a model in the loop. Therapist responses are also capped at 200 tokens,
and mean turn length roughly tripled over training in both runs, so many endpoint turns end
mid-sentence (Appendix~\ref{app:example}); the cap is identical across runs and the truncated turn
is what the patient simulator and both judges saw, so it does not bias the contrast, but it is
part of what every score in this paper measures.

\paragraph{Instruments and the utterance coder.}
\QQ\ is both the optimisation target and a headline metric, so ``\QQ\ improved'' is partly
circular: the six instruments outside the reward and the held-out judge are the checks on it
(\S\ref{sec:setup}), and Appendix~\ref{app:saturation} shows how the main judge's range narrows at
the $\Kf$ run's late checkpoints. The MITI-style coder is the least dependable instrument for
comparing the runs, and we do not lean on its differences alone. The utterance coder of
\S\ref{sec:therapist}--\ref{sec:patient} assigns one dominant code per utterance, so a
900-character turn that reflects, informs and ends in a question is coded once; its patient-side
codes reproduce the session-level change-talk counts of the same judge (pooled $\rho$
$0.88$--$0.94$ across the 21 states), but its therapist-side codes reproduce the MITI behaviour
counts only partly (Appendix~\ref{app:coder}), because the two instruments count different
things. Its most consequential therapist-side judgement, praise against affirmation, is also the
one the two judges apply differently to the $\Kf$ policy (\S\ref{sec:therapist}), and no
human coder has checked either; the sequential claims of \S\ref{sec:patient} are the claims a
human-coded sample should test first.
